\chapter{Capstone: Extending and Shipping an Assembly Game}
\label{ch:asmcap}

A program that runs only when embedded in the firmware is a lab exercise; a
program packaged as a cartridge is a game you can hand to someone. This closing
chapter of the assembly track does two things. First, it extends the Pong of
Chapter~\ref{ch:asm3} into a fuller game by reading, not rewriting, the
platform's richer examples. Second, it ships your work: it builds your source
into a \fname{.prg32} cartridge and deploys it to both QEMU and the physical
board without reflashing the firmware \citep{prg32tutorial,prg32cartfmt}.

\section{Reading is a skill: Breakout as Pong with bricks}

You will spend far more of your career reading code than writing it, and the
platform's example games are designed to be read. Once you have written Pong, the
Breakout example is legible at a glance: it is Pong with an extra piece of
state---a row of bricks---and one extra rule, removing a brick when the ball hits
it \citep{prg32repo}. The Breakout source even stores its bricks as a small bitmask,
the same bitmask idea you learned for the input in Chapter~\ref{ch:asm2}, now
repurposed so that each set bit means ``this brick is still present.''

\begin{prgconcept}
A bitmask is not only for input. Breakout packs a row of bricks into the bits of a
single word: bit $i$ set means brick $i$ is intact, clear means it has been
knocked out. Testing a brick is \code{andi}; clearing one is a bitwise AND with
the complement of its bit. The same instruction you used to read a button now
manages a row of bricks.
\end{prgconcept}

\begin{prgtry}
Open \fname{examples/games/breakout/graphics/game.S} and find the line that
declares the brick state (it is a single \code{.word} near the top, holding a
small mask). Then find, in \code{update}, where a brick is removed. Write two
sentences: one naming the register that holds the brick mask, and one explaining
which bitwise operation clears a brick. You are now reading assembly the way a
maintainer does \citep{prg32repo}.
\end{prgtry}

\section{Extending your own game}

Before packaging, extend the Pong you built in Chapter~\ref{ch:asm3} with one new
feature of your choosing. The platform's suggested progression---one moving
object, one collision, one score counter, then sound---means the natural next
features are a score that increases on a successful bounce, a second collision
surface, or a beep on the event \citep{prg32tutorial}. Add exactly one, build it,
and confirm it on QEMU. Resist the urge to add three features at once; the whole
method of this book is that a single change is easy to verify and a triple change
is not.

A score is the gentlest extension and reinforces the load--modify--store cycle
from Chapter~\ref{ch:asm1}. In \code{update}, when the collision test succeeds,
read the score word, add one, and store it back. In \code{draw}, render it with
\code{prg32\_gfx\_text8}, whose arguments are the position, the string, and a
foreground and background colour. Because turning a number into a string is
fiddly in pure assembly, a good first version simply draws a fixed label such as
\code{"SCORE"} and proves the text path works; printing the digits can come later
with the console's \code{prg32\_console\_hex32} helper, which displays a value in
hexadecimal without any string conversion on your part \citep{prg32repo}.

\begin{prgwarn}
Increment the score on the \emph{event}, not every frame. The collision test is
true for as long as the rectangles overlap, which may be several frames, so a
naive ``add one whenever overlapping'' counts a single hit many times. The robust
pattern is to act on the \emph{transition} into overlap; a simple version remembers
whether the ball was overlapping last frame and scores only when it newly
overlaps. This is the same ``act on the event'' discipline as the beep warning in
Chapter~\ref{ch:platform}.
\end{prgwarn}

\section{Building a cartridge}

When your firmware has been flashed once and your game source is ready, the
\fname{prg32\_game.py} tool compiles and packages it into a \fname{.prg32}
cartridge. The tool needs your source file, the firmware ELF so it can resolve
the framework symbols your game calls, the entry-point prefix that names your
three functions, a human-readable name, and an output path \citep{prg32tutorial}.

\begin{lstlisting}[style=bash,caption={Building an assembly cartridge against the board firmware. The \texttt{--entry-prefix} must match your \texttt{<prefix>\_init/update/draw} symbols.}]
python3 tools/prg32_game.py build \
  examples/games/pong/graphics/game.S \
  --firmware-elf build-esp32c6/PRG32.elf \
  --entry-prefix pong_graphics \
  --name pong \
  --out build-esp32c6/pong.prg32
\end{lstlisting}

\begin{prgconcept}
A cartridge bundles your three compiled functions with a small header that the
firmware reads to find them. The \code{--entry-prefix} you pass is how the tool
knows which symbols are your \code{init}, \code{update}, and \code{draw}; it must
exactly match the prefix in your \code{.global} directives. A mismatch here is a
common build-time error, and the message names the symbol it could not find.
\end{prgconcept}

\section{Deploying to both targets}

The same cartridge runs on both targets; only the delivery differs
\citep{prg32qemu}.

On the \textbf{physical board}, the resident firmware starts its own Wi-Fi access
point named \code{PRG32}. You connect your computer to that network and upload the
cartridge over HTTP; the firmware stores it and runs it, no reflashing required.

\begin{lstlisting}[style=bash,caption={Uploading a cartridge to the board over its Wi-Fi access point.}]
python3 tools/prg32_game.py upload build-esp32c6/pong.prg32 --url http://192.168.4.1
\end{lstlisting}

On \textbf{QEMU}, there is no Wi-Fi to upload over, so you build the cartridge
against the QEMU firmware ELF and stage it directly into the emulator's flash
image. If the flash image does not yet exist, run QEMU once so ESP-IDF creates it,
quit, then stage \citep{prg32qemu}.

\begin{lstlisting}[style=bash,caption={Building and staging a cartridge for the emulator. Note the QEMU firmware ELF and the \texttt{upload-qemu} subcommand.}]
python3 tools/prg32_game.py build \
  examples/games/pong/graphics/game.S \
  --firmware-elf build-qemu/PRG32.elf \
  --entry-prefix pong_graphics \
  --name pong \
  --out build-qemu/pong.prg32
python3 tools/prg32_game.py upload-qemu build-qemu/pong.prg32 --flash build-qemu/flash_image.bin
\end{lstlisting}

\begin{prgtargets}
\textbf{Both targets.} The cartridge format is identical; what changes is the
firmware ELF you build against (\fname{build-esp32c6} versus \fname{build-qemu})
and the delivery (\code{upload} over Wi-Fi versus \code{upload-qemu} into the
flash image). Your game source is byte-for-byte the same in both cases
\citep{prg32qemu}.
\end{prgtargets}

\begin{prgwarn}
Build the cartridge against the matching firmware ELF. A cartridge built against
\fname{build-esp32c6/PRG32.elf} resolves the framework symbols at the addresses
the board firmware uses; staging that same cartridge into a QEMU image built at
different addresses will not work. Build for the target you intend to run on.
\end{prgwarn}

\section{When it goes wrong}

The platform ships a self-check, \code{python3 tools/prg32\_game.py doctor}, that
verifies your toolchain prerequisites and the partition layout before you spend
time chasing a build failure \citep{prg32repo}. If a cartridge upload to QEMU
fails, the usual cause is a missing or invalid \fname{build-qemu/qemu\_flash.bin};
run QEMU once to create it, then stage again. If the board shows a black screen,
the usual cause is the wrong build directory, as warned in
Chapter~\ref{ch:platform}. These are not exotic failures; they are the everyday
friction of embedded work, and learning to read the error and apply the known fix
is itself part of the curriculum.

\begin{prgcheck}
The assembly track is complete. You can read a register-level game written by
someone else; you can extend a game by one well-tested feature; and you can
package your work as a cartridge and deploy it to both the emulator and real
hardware. You have, end to end, written and shipped a native \riscv program.
Appendix~\ref{app:worked} contains the full annotated source of the games
referenced here, and Appendix~\ref{app:api} is your reference for every framework
call you have used.
\end{prgcheck}
